\section{Introduction}

Gleason's theorem is one of the structural pillars behind the density-matrix form of quantum probability. It is also a theorem that beginners often find forbidding. Even so-called elementary proofs still demand a fair amount of geometric or analytic maturity \citep{cooke1985}. The goal of this note is not to compete with the original theorem or with the best existing proofs. The goal is more modest and, we think, more teachable.

We give a short route to a \emph{finite-dimensional real Gleason-type theorem}. The theorem concerns additive rules on the effect space
\[
\Eff_d=\{E \in \mathbb{R}^{d\times d}_{\mathrm{sym}} : 0 \leq E \leq I\}.
\]
Here $0 \leq E \leq I$ means that $E$ is positive semidefinite and all its eigenvalues are at most $1$. A rule $f:\Eff_d \to [0,1]$ is assumed to satisfy $f(I)=1$ and
\[
f(E+F)=f(E)+f(F)
\]
whenever $E$, $F$, and $E+F$ all lie in the same effect space. The conclusion is that $f(E)=\mathrm{tr}(\rho E)$ for a unique positive semidefinite trace-one matrix $\rho$.

This statement is not Gleason's original theorem. The assumption is stronger: additivity is required on the whole effect space, not merely on orthogonal projections. In that sense the note sits closer to the finite-dimensional effect-space viewpoint emphasized by \citet{busch2003} and discussed by \citet{wright2019}. What is new here is not the theorem itself, but the route: a proof staircase that begins with a one-variable lemma and keeps the ideas concrete all the way through.

\paragraph{Contributions.}
This note makes three contributions.
\begin{enumerate}[leftmargin=1.4em]
  \item It gives a bounded-additivity warmup theorem on $[0,1]$ that can be explained without physics.
  \item It turns that warmup into a sphere corollary that is accessible with only squared lengths and orthonormal triples.
  \item It uses the same bounded-additivity idea to prove the finite-dimensional real effect-space theorem in explicit matrix language.
\end{enumerate}

The computational artifacts in this package are intentionally secondary. They do not prove the theorem. They simply generate real reproducible checks and paper-side assets so the package meets the repo's standard for a complete final paper.
